\documentclass[aspectratio=169]{beamer}
\usetheme{Bruno}
\usepackage{amsmath}
\usepackage{siunitx}
\usepackage[american,RPvoltages]{circuitikz}
\ctikzset{capacitors/scale=0.7}
\usepackage{tabularx}
\newcolumntype{C}{>{\centering\arraybackslash}X}
\usepackage{tabu}
\usepackage[spanish]{babel}
\usepackage{booktabs}
\usepackage{pgfplots}
\usepgfplotslibrary{units, fillbetween} 
\pgfplotsset{compat=1.16}
\title{Electricidad I \\ \emph{Funciones de singularidad}}

\usetikzlibrary{arrows, arrows.meta}
\author{
    Juan J. Rojas
}

\institute{Instituto Tecnológico de Costa Rica}
\date{\today}
\background{fig/background.jpg}
\renewcommand\tabularxcolumn[1]{m{#1}}% for vertical centering text in X column

\begin{document}

\maketitle
\sisetup{unit-math-rm=\mathrm,math-rm=\mathrm} % change sinitx font


\begin{frame}{Funciones singulares}
    \begin{columns}
    \begin{column}{0.6\textwidth}
    \begin{itemize}
        \item La función, su derivada o ambas son discontinuas.
        \item Permiten modelar matemáticamente operaciones de conmutación de interruptores, relés y fuentes.
    \end{itemize}
    \footnotesize
    \vspace{0.2cm}
    \begin{tabularx}{\textwidth}{X l}
        \toprule
        Función & Nombre\\
        \midrule
        $u(t-t_0)$ & escalón unitario o escalón de Heaviside \\[5pt]
        $\delta(t-t_0)$ & impulso unitario o delta de Dirac \\[5pt]
        $r(t-t_0)$ & rampa unitaria  \\[5pt]
        \bottomrule
    \end{tabularx}
    \normalsize 
    \end{column}
    \begin{column}{0.4\textwidth}
    \centering
        \begin{tikzpicture}[scale=0.8]
            \begin{axis}[%
                width= 0.9\linewidth,
                height= 0.5\linewidth,
                axis lines=center,
                xlabel=$t$,
                ylabel=$u(t-t_0)$,
                x label style={anchor=west},
                y label style={anchor=south},
                ytick={1},
                ymax=1.5, % or enlarge y limits=upper
                xticklabels={$t_0$},
                xtick={2},
                xmin=-1,
                axis on top
                ]
            \addplot+[no marks, thick] coordinates {(-1,0) (2,0) (2,1) (4,1)};
            \end{axis}
        \end{tikzpicture}
        \begin{tikzpicture}[scale=0.8]
            \begin{axis}[%
                width= 0.9\linewidth,
                height= 0.5\linewidth,
                axis lines=center,
                xlabel=$t$,
                ylabel=$\delta(t-t_0)$,
                x label style={anchor=west},
                y label style={anchor=south},
                ytick=false,
                ymax=1.5, % or enlarge y limits=upper
                xticklabels={$t_0$},
                xtick={2},
                xmin=-1,
                axis on top
                ]
            \addplot+[no marks, thick] coordinates {(-1,0) (2,0) (2,1) (2,0) (4,0)};
            \draw[->,blue,thick](2,1)--(2,1.3);
            \end{axis}
        \end{tikzpicture}
        \begin{tikzpicture}[scale=0.8]
            \begin{axis}[%
                width= 0.9\linewidth,
                height= 0.5\linewidth,
                axis lines=center,
                xlabel=$t$,
                ylabel=$r(t-t_0)$,
                x label style={anchor=west},
                y label style={anchor=south},
                ytick={1},
                ymax=1.5, % or enlarge y limits=upper
                xticklabels={$t_0$,$t_0+1$},
                xtick={1,3},
                xmin=-1,
                xmax=4,
                axis on top
                ]
            \addplot+[no marks, thick] coordinates {(-1,0) (1,0) (3,1) (5,2)};
            \draw[dashed, gray] (0,1) -- (3,1) -- (3,0);
            \end{axis}
        \end{tikzpicture}
    \end{column}
\end{columns}
\end{frame}

\begin{frame}{Función impulso unitario}
    \begin{columns}
    \begin{column}{0.5\textwidth}
        \centering
        \begin{tikzpicture}[scale=1]
            \begin{axis}[%
                width= 0.9\linewidth,
                height= 0.5\linewidth,
                axis lines=center,
                xlabel=$t$,
                ylabel=$u(t-t_0)$,
                x label style={anchor=west},
                y label style={anchor=south},
                ytick={1},
                ymax=1.5, % or enlarge y limits=upper
                xticklabels={$t_0$},
                xtick={2},
                xmin=-1,
                axis on top
                ]
            \addplot+[no marks, thick] coordinates {(-1,0) (2,0) (2,1) (4,1)};
            \end{axis}
        \end{tikzpicture}
    \end{column}
    \centering
    \begin{column}{0.5\textwidth}
        \begin{equation*}
            u(t-t_0) =
            \begin{cases}
                0, & t > t_0\\
                1, & t < t_0
            \end{cases}
        \end{equation*}
        \vspace{0.5cm}
        \begin{equation*}
            \frac{d}{dt} u(t-t_0) = \delta(t-t_0)
        \end{equation*}
        \vspace{0.5cm}
        \begin{equation*}
            \int_{-\infty}^{t} u(t-t_0)dt = r(t-t_0)
        \end{equation*}
    \end{column}
\end{columns}
\end{frame}

\begin{frame}{Función escalón unitario}
    \begin{columns}
    \begin{column}{0.5\textwidth}
        \centering
        \begin{tikzpicture}[scale=1]
            \begin{axis}[%
                width= 0.9\linewidth,
                height= 0.5\linewidth,
                axis lines=center,
                xlabel=$t$,
                ylabel=$\delta(t-t_0)$,
                x label style={anchor=west},
                y label style={anchor=south},
                ytick=false,
                ymax=1.5, % or enlarge y limits=upper
                xticklabels={$t_0$},
                xtick={2},
                xmin=-1,
                axis on top
                ]
            \addplot+[no marks, thick] coordinates {(-1,0) (2,0) (2,1) (2,0) (4,0)};
            \draw[->,thick,blue](2,1)--(2,1.3);
            \end{axis}
        \end{tikzpicture}
    \end{column}
    \centering
    \begin{column}{0.5\textwidth}
        \begin{equation*}
            \delta(t-t_0) =
            \begin{cases}
                0, & t \ne t_0\\
                \infty, & t = t_0
            \end{cases}
        \end{equation*}
        \vspace{0.5cm}
        \begin{equation*}
            \int_{t_0^-}^{t_0^+} \delta(t-t_0)dt = 1
        \end{equation*}
    \end{column}
\end{columns}
\end{frame}

\begin{frame}{Resumen}
    \begin{center}
        \begin{tabularx}{14cm}{C C C}
        \toprule
        Variable & Ecuación $v-i$ & Ecuación $i-v$ \\
        \midrule
       Capacitor & $v=\frac{1}{C}\int_{t_0}^{t} i(t)dt + v(t_0) $ & $i = C\,\frac{dv}{dt}$ \\[5pt]
        Inductor & $v = L\,\frac{di}{dt}$ & $i=\frac{1}{L}\int_{t_0}^{t} v(t)dt + i(t_0) $\\[5pt]
        \bottomrule
        \end{tabularx}   
    \end{center}
\end{frame}


% \begin{frame}{Referencias}

% \bibliographystyle{ieeetr}

% \bibliography{referencias}

% \end{frame}

\end{document}